% TT14200C
% Fraktur
% 14.0
% 2013-11-30
\label{m:fraktur}
\EE{Taktik}: \Speech{\Ca{Zeitkritisch. Evtl. vitale
Bedrohung v.\,a. durch Blutverlust.}
Erkennen
von großen Blutverlusten, Blutstillung,  Steriles
Abdecken, Einschätzen der vitalen Bedrohung, ggfs
Schmerztherapie. 
Bei \E{vitaler Bedrohung}: Zeitkritisch,
zügiger Transport. 
Bei \E{stabilen Patienten}: Schienung}




\begin{itemize}
  \item  \EE{Blutstillung}: Z.\,B. mit Beckengurt
  (Beckenfraktur)
  \item Vitale Bedrohung einschätzen, ggfs. \TxMassVitMK.
  Ggfs. mit hypovolämischen Schock rechnen, Schockbekämpfung
  \item  \EE{Wundversorgung}: Bei offenen Frakturen ist
  Keimfreiheit be\-sonders wichtig, da  die Gefahr einer lebenslangen
  Knochen\-marks\-ent\-zündung \Z{(Osteomyelitis)} besteht!
  Wenn möglich soll eine \E{Wundfolie} (z.\,B.
  \index{OpSite}\T{OpSite}\texttrademark ) verwendet werden.
  \item Kleidung und Schmuck von den betroffenen Gliedmaßen
  entfernen (Schwellung!)
  \item Ggfs. Schmerztherapie (veranlassen)
  \item Ruhigstellung/Schienung bei stabilen Patienten
  \item Beurteilung von Durchblutung, Motorik und
  Sensibilität (DMS) an der betroffenen Extremität und an
  der Gegenseite
\end{itemize}